\documentclass[a4paper]{article}
% Español (México) con XeLaTeX: fontspec para fuentes Unicode, polyglossia
% para la división silábica y los nombres del documento en la variante mexicana.
\usepackage{fontspec}
\usepackage{polyglossia}
\setdefaultlanguage[variant=mexican]{spanish}
% Los nombres chinos van como «pinyin (original)»: xeCJK da a los caracteres
% Han su propia fuente; sin ella XeLaTeX los omite con un aviso.
\usepackage{xeCJK}
\setCJKmainfont{FandolSong-Regular.otf}
% polyglossia no traduce los nombres de listings.
\AtBeginDocument{\providecommand{\lstlistingname}{}\renewcommand{\lstlistingname}{Listado}%
  \providecommand{\lstlistlistingname}{}\renewcommand{\lstlistlistingname}{Índice de listados}}
\usepackage{amsmath, amssymb}
\usepackage{graphicx}
\usepackage[margin=2.5cm]{geometry}
\usepackage[most]{tcolorbox}
\usepackage{etoolbox}
\usepackage{listings}
\usepackage{hyperref}
\usepackage{booktabs}
\usepackage{float}
\newtcolorbox{knowledgebox}[1]{enhanced, colback=blue!5!white, colframe=blue!75!black, colbacktitle=blue!75!black, coltitle=white, fonttitle=\bfseries, title={#1}, attach boxed title to top left={yshift=-2mm, xshift=2mm}, boxrule=1pt, sharp corners}
\newtcolorbox{importantbox}[1]{enhanced, colback=yellow!10!white, colframe=yellow!80!black, colbacktitle=yellow!80!black, coltitle=black, fonttitle=\bfseries, title={#1}, sharp corners}
\newtcolorbox{warningbox}[1]{enhanced, colback=red!5!white, colframe=red!75!black, colbacktitle=red!75!black, coltitle=white, fonttitle=\bfseries, title={#1}, sharp corners}
\lstset{language=Python, basicstyle=\ttfamily\small, keywordstyle=\color{blue}, stringstyle=\color{red!60!black}, commentstyle=\color{green!60!black}, breaklines=true, frame=single, numbers=left, numberstyle=\tiny\color{gray}, captionpos=b, extendedchars=true}
\newcommand{\notetitle}{Irreversibilidad de codificación y decodificación del tokenizer y colapso del entrenamiento}
\newcommand{\noteauthors}{Elaboradas a partir de materiales públicos del curso}
\newcommand{\notedate}{2025}
\newcommand{\videochannel}{Wudaokou Nash (五道口纳什)}
\newcommand{\videopublishdate}{2025}
\newcommand{\videoduration}{aproximadamente 25 minutos}
\newcommand{\videourl}{}
\newcommand{\videocoverpath}{cover.jpg}
\newcommand{\repourl}{https://github.com/hqhq1025/ai-course-notes}

% Footer with repo info
\usepackage{fancyhdr}
\pagestyle{fancy}
\fancyhf{}
\fancyfoot[C]{\thepage}
\fancyfoot[R]{\footnotesize\color{gray}\href{\repourl}{AI Course Notes}}
\renewcommand{\headrulewidth}{0pt}
\renewcommand{\footrulewidth}{0pt}

\begin{document}
\begin{titlepage}\centering
{\Large Notas del curso\par}\vspace{1.2cm}{\huge\bfseries \notetitle\par}\vspace{0.8cm}{\large \noteauthors\par}\vspace{0.3cm}{\large \notedate\par}\vspace{1.2cm}
\ifdefempty{\videocoverpath}{}{\includegraphics[width=0.82\textwidth,height=0.45\textheight,keepaspectratio]{\videocoverpath}\par}
\vfill
\begin{tcolorbox}[width=0.9\textwidth, colback=black!2!white, colframe=black!60, sharp corners]
\textbf{Autor o canal del video}: \videochannel\par\textbf{Fecha de publicación}: \videopublishdate\par\textbf{Duración del video}: \videoduration\par
\ifdefempty{\videourl}{}{\textbf{Enlace del video}: \href{\videourl}{\nolinkurl{\videourl}}\par}\par
\textbf{Página del proyecto}: \href{\repourl}{\nolinkurl{\repourl}}\par
\end{tcolorbox}
\end{titlepage}
\tableofcontents\newpage

\section{Introducción}
En esta entrega se presenta un issue conocido al hacer RL post-training con veRL: el colapso del entrenamiento causado porque la codificación y decodificación del tokenizer no son reversibles.

\begin{warningbox}{Problema central}
La codificación y decodificación del tokenizer, en muchos casos, \textbf{no es reversible}:
\begin{itemize}
  \item varias listas distintas de token IDs, al decodificarse, producen el mismo texto
  \item pero al codificar ese mismo texto solo se obtiene una única lista de token IDs
  \item es decir: decode $\to$ encode $\neq$ token IDs originales
\end{itemize}
\end{warningbox}

\section{Origen del problema}
\subsection{Relación de muchos a uno en la codificación y decodificación}
Codificar y luego decodificar un texto produce un resultado determinista. Pero decodificar y luego codificar los IDs de tokens puede producir IDs de tokens distintos. Esto se debe a que varias combinaciones de IDs de tokens pueden decodificarse al mismo texto.

\subsection{Impacto en el entrenamiento de RL}
En el escenario de llamadas a herramientas de múltiples turnos:
\begin{enumerate}
  \item El modelo genera IDs de tokens (por ejemplo, en el formato de la llamada a la herramienta)
  \item La herramienta devuelve un resultado en texto
  \item El resultado se codifica de vuelta a IDs de tokens y se concatena al contexto
  \item Si el resultado de la codificación difiere de los IDs originales, la distribución de entrenamiento se sesga
  \item En casos graves, esto provoca que el entrenamiento colapse
\end{enumerate}

\subsection{El caso de Qwen2.5-7B}
Se muestra un caso concreto de inconsistencia en los IDs de tokens, que ilustra lo generalizado del problema.

\subsection{Resumen de la sección}
La irreversibilidad de la codificación y decodificación del tokenizer es un riesgo latente para el entrenamiento de RL de múltiples turnos, y requiere atención especial en el procesamiento de datos.

\newpage
\section{Línea de defensa de ingeniería: garantizar la ida y vuelta consistente entre tokens y texto}
Esta clase revela un detalle muy propio de la ingeniería, pero suficiente para arruinar el entrenamiento: si la llamada a una herramienta (tool call) o el resultado que esta devuelve cambia durante la ida y vuelta de decode/encode, el contexto de entrenamiento se desvía silenciosamente de la trayectoria original. En el RL de múltiples turnos, esta desviación se amplifica aún más en cada turno.

\begin{warningbox}{Las verificaciones que más vale la pena establecer}
\begin{itemize}
  \item Hacer pruebas de consistencia de encode/decode sobre las cadenas de formato clave
  \item Registrar en el log el resultado de recodificar el texto que devuelven las herramientas
  \item Hacer pruebas de regresión por separado para los tokenizers de alto riesgo
\end{itemize}
\end{warningbox}

\subsection{Resumen de la sección}
Los detalles del tokenizer no son «de tan bajo nivel que se puedan ignorar». En el RL de múltiples turnos, determinan directamente si el contexto es estable y reproducible.

\newpage
\section{Síntesis y ampliación}
\begin{enumerate}
  \item Tokenizer encode-decode no siempre es reversible
  \item Esto provoca un sesgo de distribución en el entrenamiento de Multi-turn RL
  \item Solución: garantizar un manejo consistente de los token IDs
\end{enumerate}
\end{document}

